\section{Background}

\subsection{LLVM IR and SSA}

LLVM IR is a typed intermediate representation in static single assignment form.
Every value is defined exactly once, and a use always refers to a single
definition, so the data flow of a function is explicit in its structure. Control
flow is a graph of basic blocks, each a straight line sequence of instructions
ending in a terminator (a branch, a return, and so on). Where control flow
merges, a phi node selects a value based on which predecessor block was taken.
The full semantics are specified in the language reference \cite{llvm_langref}.

This form is what makes middle end optimization tractable. Because each value has
one definition, an analysis can reason about a computation without tracking
assignments, and a transform can replace all uses of a value in one operation.
The transforms in this project rely on exactly that: replacing a redundant
computation, or an identity such as an add of zero, is a single replace all uses
followed by an erase.

\subsection{The new pass manager}

The new pass manager is LLVM's current framework for running middle end work
\cite{llvm_writing_pass}. A transform is a small type with a \texttt{run} method
returning a \texttt{PreservedAnalyses} set that states precisely which cached
analyses survive the change. An analysis is a type that computes a result and is
cached by an analysis manager, which invalidates the cache when a transform does
not preserve it. Getting that contract right is central: a transform that claims
to preserve an analysis it actually invalidated leaves later passes reading stale
facts.

A plugin registers its passes through \texttt{PassBuilder} callbacks and exposes
them by name, so a stock \texttt{opt} can run them from a pipeline string such as
\texttt{-passes="my-canon,my-simplify,my-dce"} with no patched binary
\cite{llvm_project}. This project is such a plugin, built against a distribution
LLVM with no source build of the compiler itself.
